% !TeX root = ../tesi.tex

%definisco il layout dell'abstract
\def\changemargin#1#2{\list{}{\rightmargin#2\leftmargin#1}\item[]}
\let\endchangemargin=\endlist

%Genero l'ambiente per l'abstract
\newcommand\summaryname{Abstract}
\newenvironment{Abstract}%
    {\begin{center}%
    \bfseries{\summaryname} \end{center}}

\begin{Abstract}
    \begin{changemargin}{1cm}{1cm}
       This thesis presents an embedded multimodal computer-vision framework for quality monitoring in a real industrial bread-production process, designed to operate in an edge environment without network connectivity. The objective is to introduce minimally invasive technological tools into a real, weakly structured industrial environment, enabling both automated quality inspection and statistical analysis of the production process. The system was developed and validated in the \emph{SMACT Competence Center}, as part of a collaborative project led by FlexSight Srl in the context of the Vis.IO project, developed for IRINOX S.p.A. under the IRISS framework (Ref. COR 22612178, CUP H19J24001010004 - funded by NextGenerationEU, M4C2I2.3). The framework was deployed at two acquisition sites characterized by different sensing modalities and inspection objectives. At Site~1, an NVIDIA Jetson Nano DK coupled with an Azure Kinect acquires RGB-D data at the output of the dough-processing stage. At Site~2, a Raspberry Pi 5 equipped with a Raspberry Pi Camera Module~3 Wide acquires RGB images of baked bread during the final handling stage.

        In both sites, product localization is performed through a DEIMv2-based object-detection pipeline trained from automatically labelled data. The downstream analysis then differs according to the production stage. At Site~1, the detected dough balls are segmented and converted into local 3D point clouds in order to estimate product volume through different geometric reconstruction strategies. At Site~2, the detected bread loaves are cropped and analyzed with an anomaly-detection model to assess the quality of the final baked product.

        The experimental results confirm the effectiveness of the proposed framework in both scenarios. For dough-ball inspection, the rotated convex-hull method yields the most accurate volume estimates, showing the closest agreement with the expected reference interval among the evaluated methods. For baked bread inspection, the anomaly-detection branch achieved a near-perfect image-level AUROC of 0.999958 under the adopted evaluation protocol, while also yielding very low anomaly scores on normal samples and meaningful responses on anomalous examples. The corresponding anomaly maps localize the most relevant irregular regions.
        \newline

        Overall, the thesis shows that computer vision, deep learning, and 3D sensing can be effectively integrated into minimally invasive embedded systems to support data-driven quality monitoring in real industrial food-production environments.
        % changemargin %levalo per far tornare i margini alla dimensione corretta
    \end{changemargin}
\end{Abstract}
